\documentclass[12pt]{article}
% Shared preamble and text for the data-request letters.
% Replace the bracketed fields before submitting any letter.

\usepackage[a4paper,top=2.4cm,bottom=2.4cm,left=2.8cm,right=2.4cm]{geometry}
\usepackage[utf8]{inputenc}
\usepackage[T1]{fontenc}
\usepackage{lmodern}
\usepackage{microtype}
\usepackage{setspace}
\usepackage{array}
\usepackage{booktabs}
\usepackage{enumitem}
\usepackage{tabularx}
\usepackage{xcolor}
\usepackage[hidelinks]{hyperref}

\setstretch{1.12}
\setlength{\parindent}{0pt}
\setlength{\parskip}{6pt}
\setlist[itemize]{leftmargin=1.5em,itemsep=2pt,topsep=2pt}
\setlist[enumerate]{leftmargin=1.7em,itemsep=2pt,topsep=2pt}
\newcolumntype{Y}{>{\raggedright\arraybackslash}X}
\newcolumntype{P}[1]{>{\raggedright\arraybackslash}p{#1}}

\newcommand{\NamaMahasiswa}{[NAMA LENGKAP MAHASISWA]}
\newcommand{\NIMMahasiswa}{[NIM]}
\newcommand{\NamaPembimbing}{[NAMA DAN GELAR DOSEN PEMBIMBING]}
\newcommand{\JudulSementara}{Diagnosis \textit{Borrowed Size} dan \textit{Agglomeration Shadow} antara Jakarta dan Pusat-pusat Sekunder di Kawasan Aglomerasi Jakarta}

\newcommand{\KopSurat}{%
  \begin{center}
    \fbox{%
      \begin{minipage}{0.88\textwidth}
        \centering
        \textbf{KOP SURAT RESMI FAKULTAS TEKNIK UNIVERSITAS GADJAH MADA}\\
        \small Ganti blok ini dengan kop resmi sebelum surat dikirim.
      \end{minipage}}
  \end{center}
  \vspace{4pt}
}

\newcommand{\IdentitasMahasiswa}{%
  \begin{tabularx}{\textwidth}{@{}P{0.30\textwidth}Y@{}}
    Nama & \NamaMahasiswa \\
    NIM & \NIMMahasiswa \\
    Program studi & Sarjana Perencanaan Wilayah dan Kota \\
    Departemen dan fakultas & Departemen Teknik Arsitektur dan Perencanaan, Fakultas Teknik, Universitas Gadjah Mada \\
    Dosen pembimbing & \NamaPembimbing
  \end{tabularx}
}

\newcommand{\PernyataanPenggunaanData}{%
Data yang dimohon akan digunakan untuk kepentingan pendidikan dan penelitian nonkomersial. Data tidak akan dijual, dialihkan, atau disebarluaskan dalam bentuk mentah. Hasil yang dipublikasikan akan berupa statistik agregat, peta turunan, atau keluaran lain yang diizinkan oleh pemilik data. Pengolahan dan penyimpanan data akan mengikuti ketentuan kerahasiaan, lisensi, pembatasan publikasi, dan perjanjian penggunaan data yang berlaku.
}

\newcommand{\PernyataanTanpaBiaya}{%
Apabila permintaan ini tidak dapat dipenuhi tanpa biaya, mohon kami diberi tahu terlebih dahulu. Kami tidak menyetujui pembuatan tabulasi atau pengadaan data berbayar tanpa persetujuan baru dari mahasiswa dan dosen pembimbing.
}

\newcommand{\AbstraksiPenelitian}{%
\textbf{Latar belakang.} Kegiatan ekonomi, perjalanan, layanan, dan perumahan di sekitar Jakarta berlangsung melampaui batas administrasi DKI Jakarta. Hubungan tersebut tidak otomatis menunjukkan bahwa pusat-pusat sekunder telah memperkuat fungsi lokalnya. Suatu wilayah dapat memperoleh akses ke pasar metropolitan, tetapi pada saat yang sama semakin bergantung pada Jakarta untuk pekerjaan, layanan, atau fungsi bernilai tinggi.

\textbf{Tujuan.} Penelitian ini mengkaji apakah integrasi dengan Jakarta berjalan bersama penguatan fungsi lokal, suatu pola yang konsisten dengan \textit{borrowed size}, atau bersama defisit fungsi dan ketergantungan yang asimetris, suatu pola yang konsisten dengan \textit{agglomeration shadow}. Penelitian tidak bertujuan mengusulkan perluasan wilayah DKI Jakarta, mengubah batas pemerintahan, atau membuktikan hubungan kausal hanya dari potret satu periode.

\textbf{Metode.} Kasus utama adalah sistem metropolitan Jabodetabek. Wilayah Statistik Metropolitan Indonesia digunakan sebagai sabuk pembanding. DKI Jakarta diperlakukan sebagai satu inti dengan batas administrasi tetap. Kecamatan menjadi unit utama jika resolusi data mendukung. Data pada tingkat kabupaten/kota, zona transportasi, raster, dan titik dipertahankan pada resolusi aslinya sebelum dilakukan agregasi yang dapat ditelusuri. Analisis membedakan data yang teramati, tabulasi agregat, hasil pemodelan, dan proksi.

\textbf{Cakupan wilayah dan waktu.} Baseline utama menggunakan tahun 2024. Data tahun lain digunakan sebagai konteks atau pemeriksaan perubahan, dengan tahun dan perbedaan unit ditampilkan secara terbuka. Wilayah mencakup DKI Jakarta, kabupaten/kota dalam kawasan utama Jabodetabek, pusat sekunder yang ditetapkan melalui prosedur yang terdokumentasi, serta sabuk pembanding yang diperlukan untuk menguji fungsi relatif.

\textbf{Jenis data dan variabel.} Penelitian memerlukan data penduduk, luas, kepadatan, integrasi dan perjalanan, fungsi layanan lokal, pekerjaan menurut lokasi tempat kerja dan sektor jika tersedia, aksesibilitas jaringan, serta karakteristik pusat sekunder. Podes digunakan untuk mengukur kehadiran dan akses layanan, bukan sebagai ukuran langsung produktivitas atau jumlah pekerjaan. Data komuter pada tingkat kabupaten/kota tidak akan dipaksa menjadi estimasi kecamatan. Data OD yang dimodelkan akan diberi label sintetis dan diuji terhadap kendala yang tersedia.

\textbf{Rencana hasil.} Hasil penelitian berupa profil integrasi dan fungsi lokal, tipologi kondisi relasional, peta gradien, uji sensitivitas bobot dan ambang, serta penjelajah skenario berbasis WebGL. Penjelajah tersebut menghitung ulang dan menampilkan hasil secara interaktif. Ia bukan sumber data waktu nyata, ramalan kebijakan, atau simulasi perubahan kewenangan pemerintahan.
}

\newcommand{\LampiranAbstraksi}{%
  \clearpage
  \section*{Lampiran 2. Abstraksi penggunaan data}
  \AbstraksiPenelitian
}

\newcommand{\TandaTanganDTAP}{%
  Hormat kami,\\[2.2cm]
  \parbox[t]{0.85\textwidth}{[NAMA PENANDATANGAN]}\\
  \parbox[t]{0.85\textwidth}{[JABATAN PENANDATANGAN SESUAI ALUR PERSURATAN DTAP]}
}

\newcommand{\TandaTanganBPS}{%
  Hormat kami,\\[2.2cm]
  \parbox[t]{0.85\textwidth}{[NAMA PENANDATANGAN]}\\
  \parbox[t]{0.85\textwidth}{[DEKAN FAKULTAS TEKNIK ATAU PEJABAT SETINGKAT YANG DITERIMA BPS]}
}


% The BPS leaflet states that the letter and the abstraction are uploaded
% together as one PDF. Compile this file as one submission document.

\begin{document}
\KopSurat

\begin{flushright}
Yogyakarta, [TANGGAL SURAT]
\end{flushright}

\begin{tabularx}{\textwidth}{@{}P{0.20\textwidth}Y@{}}
  Nomor & [DIISI OLEH UNIT PERSURATAN] \\
  Lampiran & Dua lampiran \\
  Hal & Permohonan layanan data mikro melalui mekanisme tarif Rp0,00
\end{tabularx}

Yth. [UNIT PENERIMA SESUAI KONFIRMASI LAYANAN SILASTIK]\\
Badan Pusat Statistik Republik Indonesia\\
Jakarta

Dengan hormat,

Sehubungan dengan penyusunan Tugas Akhir pada Program Sarjana Perencanaan Wilayah dan Kota, Departemen Teknik Arsitektur dan Perencanaan, Fakultas Teknik, Universitas Gadjah Mada, kami mengajukan permohonan layanan data untuk mahasiswa berikut.

\IdentitasMahasiswa

Penelitian berjudul sementara ``\JudulSementara''. Penelitian ini mengkaji hubungan fungsional antara DKI Jakarta dan pusat-pusat sekunder di kawasan Jabodetabek. Data diperlukan untuk membaca fungsi layanan lokal, pola komuter, ukuran wilayah, serta pekerjaan menurut lokasi tempat kerja apabila tabulasi tersebut dapat disajikan pada unit yang sah.

Kami memohon agar produk dan permintaan pada Lampiran 1 dipertimbangkan melalui mekanisme tarif Rp0,00 bagi institusi pendidikan dalam negeri, sesuai ketentuan yang berlaku dan setelah kelayakan pemohon dikonfirmasi oleh BPS. Podes 2024 dan Survei Komuter Jabodetabek 2023 menjadi prioritas pertama. Permintaan Sakernas Agustus 2024 bersifat kondisional dan diutamakan dalam bentuk subset wilayah atau tabulasi khusus, bukan file nasional penuh. Podes 2021 hanya diperlukan apabila kapasitas layanan masih memungkinkan setelah prioritas pertama diproses.

Apabila beberapa produk tidak dapat digabung dalam satu transaksi atau satu permohonan, mohon diinformasikan produk mana yang dapat diproses terlebih dahulu dan apakah diperlukan transaksi atau surat terpisah. Apabila data mikro tidak dapat diberikan, kami bersedia menerima tabulasi agregat, unit terkecil yang memenuhi standar reliabilitas, atau penjelasan tertulis mengenai bentuk data yang dapat disediakan.

\PernyataanPenggunaanData

\PernyataanTanpaBiaya

Kami bersedia memenuhi persyaratan administrasi dan menandatangani Surat Perjanjian Penggunaan Data apabila permohonan disetujui, setelah isi perjanjian diperiksa oleh mahasiswa, dosen pembimbing, dan pejabat yang berwenang.

Demikian permohonan ini kami sampaikan. Atas perhatian dan bantuan Bapak/Ibu, kami mengucapkan terima kasih.

\vspace{8pt}
Tembusan:
\begin{enumerate}
  \item Dosen pembimbing;
  \item Ketua Program Studi Sarjana Perencanaan Wilayah dan Kota;
  \item Arsip.
\end{enumerate}

\TandaTanganBPS

\clearpage
\section*{Lampiran 1. Rincian data yang dimohon}

\renewcommand{\arraystretch}{1.25}
\begin{tabularx}{\textwidth}{@{}P{0.14\textwidth}P{0.30\textwidth}Y@{}}
\toprule
Prioritas & Produk atau permintaan & Unit, periode, dan kelengkapan yang dimohon \\
\midrule
P0 & Mikrodata Podes 2024 & Tingkat desa atau kelurahan untuk wilayah studi. Mohon disertakan kode wilayah, kamus data, kuesioner atau metadata yang dapat dibagikan, serta catatan perubahan definisi. \\
P0 & Mikrodata Survei Komuter Jabodetabek 2023 & Mohon disertakan dokumentasi desain sampel, bobot, unit analisis, dan variabel yang diizinkan untuk asal, tujuan, moda, durasi, jarak, biaya, pekerjaan, serta karakteristik responden. \\
P0 kondisional & Sakernas Agustus 2024 & Diutamakan subset Jabodetabek atau DKI, Jawa Barat, dan Banten, atau tabulasi jumlah pekerja menurut lokasi tempat kerja dan kelompok sektor luas pada unit terkecil yang layak. Mohon disertakan bobot, ukuran reliabilitas, dan aturan supresi. \\
P1 & Podes 2021 & Hanya untuk pemeriksaan perubahan fungsi 2021 sampai 2024. Mohon disertakan kamus data dan keterangan perubahan kode atau definisi jika produk ini dapat diproses tanpa menggeser P0. \\
Kondisional & Peta digital Wilkerstat dan crosswalk & Hanya jika diperlukan untuk mengharmonisasi kode desa, kecamatan, dan wilayah studi. Mohon disertakan tahun, sistem koordinat, kode, dan ketentuan penggunaan. \\
\bottomrule
\end{tabularx}

\vspace{8pt}
Jika batas layanan atau kelayakan pemohon membatasi permintaan, kami memohon urutan berikut: Podes 2024, Survei Komuter Jabodetabek 2023, Sakernas dalam bentuk subset atau tabulasi yang layak, lalu Podes 2021. Kami menerima subset wilayah Greater Jakarta dan variabel yang langsung berkaitan dengan tujuan penelitian.

\begin{itemize}
  \item Kami tidak meminta nama, alamat, NIK, nomor kontak, jejak perjalanan individu, atau pengenal usaha yang bersifat rahasia.
  \item Kami tidak meminta BPS menurunkan data kabupaten atau kota menjadi data kecamatan secara artifisial.
  \item Jika suatu produk tidak termasuk layanan tarif Rp0,00, mohon beri tahu sebelum proses berbayar dilakukan.
\end{itemize}

\LampiranAbstraksi

\end{document}
